\chapter{cna-samples and cna-examples}
\label{ch:samples-and-examples}

The two sibling repositories answer different questions.  \texttt{cna-samples} translates
Microsoft's historical XNA sample programs and exposes compatibility pressure.  \texttt{cna-examples}
is one CNA-native catalog application designed to demonstrate the current API.  Neither is a
substitute for CNA's tests, and neither repository's status prose should be read without its pinned
revision and evidence type.

\section{cna-samples: ports with provenance}

At commit \texttt{149a19f}, the repository tracks the official XNA Game Studio~4.0 collection
under the Microsoft Permissive License.  It expects CNA and sharp-runtime as sibling directories
and builds each completed sample as an independent executable.

Its plan accounts for 153 source-archive directories in three populations: 63 completed ports,
23 tracked placeholders, and 67 explicitly ignored entries.  The filesystem corroborates the
first two numbers: 86 sample directories each carry a \texttt{missing.md}, and 63 contain a sample
CMake target.  The ignored set covers tools, redundant training material, platform/service-bound
programs, and other declared non-targets rather than pretending every archive directory is a
portable game.

These are catalog statuses, not fresh run results.  ``Done'' means real port source and a build
target exist; it does not mean pixel equality, all original scenes, or all current CNA renderers.
The plan itself records the instructive SimpleAnimation correction: an entry already marked done
later revealed two CNA graphics defects when re-reviewed.  A sample is excellent integration
pressure precisely because its behavior crosses boundaries that unit tests may isolate away.

\section{Read each missing file as a dated case record}

Every tracked sample's \texttt{missing.md} explains whether it is complete, simplified, blocked,
or deliberately scoped.  These files record concrete asset substitutions, content-reader gaps,
hardware fallbacks, and framework defects found during translation.  They are better evidence
than a single top-level ratio because they name the affected program and mechanism.

They also age.  The pinned top-level plan still says XNB is unsupported, while CNA now has a real
XNB stack.  Some historical paths point to repositories or layouts that no longer exist.  A
placeholder must therefore be re-tested against current CNA before it is counted as still blocked.
The right workflow is: read the old reason, locate the current API and test, then either port the
sample or replace the reason with a new, evidenced boundary.

\section{Why official samples remain valuable}

An official sample combines API calls in the order and lifecycle Microsoft taught users to rely
on.  That makes it good at finding defects such as default-state mismatch, model winding,
component timing, asset transformation, or a renderer path that works in isolation but not in a
whole game.

It is not automatically an independent correctness oracle.  A C\texttt{++} translation can
misread the C\# source, replace assets, skip a scene, or work around a CNA gap.  The strongest use
pairs the port with the original sample or another implementation, preserves inputs, and compares
screenshots or traces.  Only two entries in the plan have direct official MonoGame sample
counterparts, so that differential opportunity is narrower than the catalog size suggests.

\section{cna-examples: one navigable CNA-native application}

At commit \texttt{fa0b49a}, \texttt{cna-examples} contains one application with 13 areas, 79
categories, and 249 registered demo screens.  A read-only run of its
\texttt{tools/check\_catalog.py} at that pin confirms every screen is defined, registered exactly
once, and documented.  Areas span framework, math, content, storage, diagnostics, input, audio,
devices, networking, media, avatars, and 2D/3D graphics.

The application offers menu navigation, substring search, direct \texttt{--demo} selection,
scripted actions, bounded frame counts, and screenshots.  This makes a screen useful both to a
human exploring the API and to a headless sweep.  Global state restoration matters because all
screens share one process; the Framework sweep explicitly checks that a resolution-changing demo
returns the back buffer to its prior size.

\section{A screenshot sweep needs semantic companions}

The repository distinguishes several verification forms:

\begin{itemize}
  \item real hardware interaction for keyboard, mouse, and much of audio;
  \item Xvfb screenshot/behavior sweeps for Framework and Media;
  \item C\texttt{++} claim programs for Math, CNJ, and XACT statements rendered on screen;
  \item self-checking 3D screens that turn readback or occlusion results into a visible verdict;
  \item graceful explanatory screens for unavailable controller, touch, sensor, camera, or dialog
        hardware.
\end{itemize}

This separation prevents a caught exception from painting a clean explanatory screen and being
mistaken for a successful API demonstration.  A screenshot checker can detect blank output or
overflow; a claim program or self-check must decide whether the described operation was correct.

\section{The pinned documentation uses obsolete renderer names}

\texttt{cna-examples} predates CNA's renderer terminology migration.  Its README and build
examples use the retired backend-named option with internal value \texttt{EASYGL}; current CNA
uses \texttt{CNA\_GRAPHICS\_RENDERER} with \texttt{OPENGLES3} or another public identity.  The repository's recorded
249-of-249 EasyGL sweep refers to the implementation family now selected publicly as an EasyGL-
backed identity.  It must not be copied as a current configure command.

The SDL-renderer sweep remains conceptually useful: its 2D-only capability causes 3D demos to
explain their unavailability rather than throw.  Yet the recorded sweep is a dated result, not a
run performed during this book audit.  Revalidation should update selector names first and retain
screenshots plus tool output.

\section{Use both repositories as triangulation}

For learning, start with \texttt{cna-examples}: one build, searchable screens, current CNA-shaped
code, and explicit evidence helpers.  For migration risk, inspect \texttt{cna-samples}: original
program provenance, per-sample scope records, and failures driven by realistic combinations.

For framework confidence, make the two disagree productively.  If a CNA-native screen passes but
an official port fails, inspect translation and lifecycle.  If both fail across one renderer but
pass another, localize the renderer.  If both pass yet differ from XNA, bring in a real-XNA or FNA
oracle.  Examples demonstrate; samples integrate; tests and independent references decide.
